% Generated draft: requires mathematical and source review.
\documentclass[11pt]{article}
\usepackage[margin=25mm]{geometry}
\usepackage{fontspec}
\setmainfont{Latin Modern Roman}
\usepackage{amsmath,amssymb,mathtools,unicode-math}
\setmathfont{Latin Modern Math}
\usepackage{enumitem,needspace,xurl,hyperref}
\setlength{\parindent}{0pt}
\setlength{\parskip}{5pt}
\setlength{\emergencystretch}{3em}
\newcommand{\N}{\mathbb{N}}
\newcommand{\Z}{\mathbb{Z}}
\newcommand{\Q}{\mathbb{Q}}
\newcommand{\R}{\mathbb{R}}
\newcommand{\C}{\mathbb{C}}
\newcommand{\F}{\mathbb{F}}
\newcommand{\A}{\mathbb{A}}
\newcommand{\PP}{\mathbb{P}}
\newcommand{\E}{\mathbb{E}}
\newcommand{\eps}{\varepsilon}
\newcommand{\dd}{\,\mathrm{d}}
\begin{document}
\section*{UnsolvedMath ID: 10900061}
\textbf{Source problem code:} catalogue entry
\subsection{Definitions and mathematical statement}
The source recalls the Gay--Kirby trisection theorem: every closed orientable smooth $4$-manifold admits a trisection, and any two trisections have common stabilizations. A trisection decomposes the manifold into three $4$-dimensional $1$-handlebodies, with pairwise intersections $3$-dimensional handlebodies and triple intersection a closed surface.

For context, a customary higher-dimensional multisection model in dimension $n=2k$ or $n=2k+1$ uses $k+1$ pieces
\[
 M=H_0\cup\cdots\cup H_k,
\]
with disjoint interiors, each $H_i$ an $n$-dimensional $1$-handlebody, and with controlled multiple intersections. The common intersection has dimension $n-k$; an intersection of $r$ pieces has dimension $n-r+1$ before the final common intersection. The additional spine/handle conditions on these intersections, smooth corner conventions, and any balancing requirements are part of a precise definition, not optional conclusions. The supplied excerpt does not state them, so this contextual description is not a substitute for the article's intended complete definition.

\textbf{Problems (Tillmann).}
\begin{enumerate}
\item Do higher-dimensional smooth manifolds always admit multisections? In a fixed manifold category and multisection convention, this asks whether every manifold in the intended class has such a decomposition. The preceding $4$-dimensional theorem concerns closed orientable smooth manifolds; the excerpt does not explicitly say whether those restrictions also govern the higher-dimensional question or whether boundary, noncompact, or nonorientable cases are included.
\item For smooth $n$-manifolds with $n\ge5$, what is the correct counterpart of uniqueness up to stabilization? Specify the stabilization operations and the appropriate equivalence relation on multisections, and determine the resulting uniqueness statement. The motivating form would give a common stabilization for any two multisections of a fixed manifold, up to the specified equivalence, but the question explicitly asks for the correct generalization rather than supplying an already defined higher-dimensional stabilization theorem.
\end{enumerate}
Both the existence question and the stabilization/uniqueness question are retained. PL existence results, if used, need an explanation of their relation to the smooth category requested here.
\subsection{Short English statement}
Do all higher-dimensional smooth manifolds admit multisections, and what stabilization operations give the correct uniqueness statement in dimensions at least five? The manifold class, full multisection definition, and equivalence convention need confirmation.
\subsection{Sources}
\begin{itemize}
\item Local UnsolvedMath registry, permanent ID 10900061, source code catalogue entry; attributed to Tillmann, no proposal year supplied.
\item Delp, Hoffoss and Manning, \emph{Problems In Groups, Geometry, and Three-Manifolds} (2015), subsection 8.7 (Tillmann), PDF page 6; \url{https://arxiv.org/pdf/1512.04620#page=6} and \url{https://arxiv.org/abs/1512.04620}. The provenance records source-TeX extraction checked against rendered PDF text and access on 2026-07-29.
\item The source attributes the $4$-dimensional existence and common-stabilization theorem to Gay--Kirby. The summary reports higher-dimensional existence progress and unresolved general uniqueness, without separate primary bibliographic citations.
\end{itemize}
\end{document}
